\centerline{\bf \Huge Текстовые задачи}

\problem{\textit{\textbf{1}}} Из пункта $A$ в пункт $B$ одновременно выехали два автомобиля. Первый проехал с постоянной скоростью весь путь. Второй проехал первую половину пути со скоростью, меньшей скорости первого на 13 км/ч, а вторую половину пути — со скоростью 78 км/ч, в результате чего прибыл в пункт В одновременно с первым автомобилем. Найдите скорость первого автомобиля, если известно, что она больше 48 км/ч. Ответ дайте в км/ч.

\problem{\textit{\textbf{2}}} Из городов $A$ и $B$ навстречу друг другу выехали мотоциклист и велосипедист. Мотоциклист приехал в $B$ на 3 часа раньше, чем велосипедист приехал в $A$, а встретились они через 48 минут после выезда. Сколько часов затратил на путь из $B$ в $A$ велосипедист?

\problem{\textit{\textbf{3}}} Первый велосипедист выехал из поселка по шоссе со скоростью 15 км/ч. Через час после него со скоростью 10 км/ч из того же поселка в том же направлении выехал второй велосипедист, а еще через час после этого — третий. Найдите скорость третьего велосипедиста, если сначала он догнал второго, а через 2 часа 20 минут после этого догнал первого. Ответ дайте в км/ч.

\problem{\textit{\textbf{4}}} Первые 190 км автомобиль ехал со скоростью 50 км/ч, следующие 180 км— со скоростью 90 км/ч, а затем 170 км — со скоростью 100 км/ч. Найдите среднюю скорость автомобиля на протяжении всего пути. Ответ дайте в км/ч.

\problem{\textit{\textbf{5}}} Поезд, двигаясь равномерно со скоростью 60 км/ч, проезжает мимо лесополосы, длина которой равна 400 метрам, за 1 минуту. Найдите длину поезда в метрах.

\problem{\textit{\textbf{6}}} Из одной точки круговой трассы, длина которой равна 14 км, одновременно в одном направлении стартовали два автомобиля. Скорость первого автомобиля равна 80 км/ч, и через 40 минут после старта он опережал второй автомобиль на один круг. Найдите скорость второго автомобиля. Ответ дайте в км/ч.

\problem{\textit{\textbf{7}}} Два гонщика участвуют в гонках. Им предстоит проехать 60 кругов по кольцевой трассе протяжённостью 3 км. Оба гонщика стартовали одновременно, а на финиш первый пришёл раньше второго на 10 минут. Чему равнялась скорость второго гонщика, если известно, что первый гонщик в первый раз обогнал второго на круг через 15 минут? Ответ дайте в км/ч.

\newpage

\hspace{3mm}

\problem{\textit{\textbf{8}}} Моторная лодка прошла против течения реки 112 км и вернулась в пункт отправления, затратив на обратный путь на 6 часов меньше. Найдите скорость течения, если скорость лодки в неподвижной воде равна 11 км/ч. Ответ дайте в км/ч.

\problem{\textit{\textbf{9}}} Весной катер идёт против течения реки в $\dfrac{5}{3}$ раза медленнее, чем по течению. Летом течение становится на 1 км/ч медленнее. Поэтому летом катер идёт против течения в $\dfrac{3}{2}$ раза медленнее, чем по течению. Найдите скорость течения весной (в км/ч).

\problem{\textit{\textbf{10}}} Первая труба пропускает на 5 литров воды в минуту меньше, чем вторая. Сколько литров воды в минуту пропускает вторая труба, если резервуар объемом 375 литров она заполняет на 10 минут быстрее, чем первая труба заполняет резервуар объемом 500 литров?

\problem{\textit{\textbf{11}}} Даша и Маша пропалывают грядку за 12 минут, а одна Маша — за 20 минут. За сколько минут пропалывает грядку одна Даша?

\problem{\textit{\textbf{12}}} Заказ по изготовлению деталей ученик токаря может выполнить за 18 часов, а токарь — за 12 часов. Ученик начал выполнять такой заказ. Через какое время после начала выполнения заказа учеником нужно начать работу токарю, чтобы в этом заказе деталей, изготовленных учеником, было в два раза больше деталей, изготовленных токарем? Ответ дайте в часах.

\problem{\textit{\textbf{13}}} Изюм получается в процессе сушки винограда. Сколько килограммов винограда потребуется для получения 20 килограммов изюма, если виноград содержит 90$\%$ воды, а изюм содержит 5$\%$ воды?

\problem{\textit{\textbf{14}}} Имеется два сплава. Первый содержит 10$\%$ никеля, второй — 30$\%$ никеля. Из этих двух сплавов получили третий сплав массой 200 кг, содержащий 25$\%$ никеля. На сколько килограммов масса первого сплава была меньше массы второго?

\problem{\textit{\textbf{15}}} Имеются два сосуда. Первый содержит 30 кг, а второй — 20 кг раствора кислоты различной концентрации. Если эти растворы смешать, то получится раствор, содержащий 68$\%$ кислоты. Если же смешать равные массы этих растворов, то получится раствор, содержащий 70$\%$ кислоты. Сколько килограммов кислоты содержится в первом сосуде?